\newpage
\phantomsection
\label{prf:ec-negation}

\begin{remark*}[Return]
\hyperref[thm:ec-negation]{Return to Theorem}
\end{remark*}

\begin{theorem*}[Negation of an Equivalence Class]
For every $[(a,b)]\in\mathbb{Q}$,
\[
-[(a,b)] = [(-a,b)] = [(a,-b)].
\]
\end{theorem*}

\begin{proof}
\textbf{Professional Standard Proof.}~\\
Let $[(a,b)]\in\mathbb{Q}$. We must show that $[(-a,b)]$ is the additive
inverse of $[(a,b)]$ in $\mathbb{Q}$, and that $[(-a,b)] = [(a,-b)]$.

\textbf{Part 1.} $[(-a,b)]$ is the additive inverse of $[(a,b)]$.

Compute the sum using the definition of addition:
\[
[(a,b)] + [(-a,b)]
= [(a\cdot b + b\cdot(-a),\; b\cdot b)]
= [(ab - ab,\; b^2)]
= [(0,\; b^2)].
\]
Since $b\neq 0$, we have $0\cdot b^2 = 0 = 0\cdot 1$, so $(0,b^2)\sim(0,1)$
and thus $[(0,b^2)] = [(0,1)] = 0$.

Therefore $[(a,b)] + [(-a,b)] = 0$. By uniqueness of additive inverses in
any group, $-[(a,b)] = [(-a,b)]$.

\textbf{Part 2.} $[(-a,b)] = [(a,-b)]$.

By the definition of $\sim$, this holds iff $(-a)(-b) = b\cdot a$. Since
$(-a)(-b) = ab = ba$, the equality holds. Hence $[(-a,b)] = [(a,-b)]$.
\end{proof}

\begin{proof}
\textbf{Detailed Learning Proof.}~\\

\textbf{Step 1.} Recall what it means to be an additive inverse.

The additive inverse of $[(a,b)]$ in $\mathbb{Q}$ is the unique element
$x\in\mathbb{Q}$ satisfying
\[
[(a,b)] + x = 0 = [(0,1)].
\]
Additive inverses are unique in any group (since $\mathbb{Q}$ under
addition is an abelian group). We will verify this identity for
$x = [(-a,b)]$.

\textbf{Step 2.} Check that $[(-a,b)]\in\mathbb{Q}$.

Since $b\neq 0$ (as $(a,b)\in W$) and $-a\in\mathbb{Z}$, the pair
$(-a,b)$ is an admissible fraction-pair, so $[(-a,b)]\in\mathbb{Q}$.

\textbf{Step 3.} Compute $[(a,b)] + [(-a,b)]$.

By the definition of rational addition,
\[
[(a,b)] + [(-a,b)]
= [(a\cdot b + b\cdot(-a),\; b\cdot b)].
\]
The numerator is
\[
a\cdot b + b\cdot(-a) = ab - ab = 0.
\]
The denominator is $b^2$. Since $b\neq 0$, we have $b^2\neq 0$, so the
result is $[(0,b^2)]$.

\textbf{Step 4.} Show that $[(0,b^2)] = [(0,1)]$.

By the definition of $\sim$, $(0,b^2)\sim(0,1)$ iff $0\cdot 1 = b^2\cdot 0$.
Both sides equal $0$, so the equivalence holds. Therefore $[(0,b^2)] = 0$.

\textbf{Step 5.} Conclude that $-[(a,b)] = [(-a,b)]$.

We have shown $[(a,b)] + [(-a,b)] = 0$. By uniqueness of additive inverses,
\[
-[(a,b)] = [(-a,b)].
\]

\textbf{Step 6.} Show that $[(-a,b)] = [(a,-b)]$.

By the definition of $\sim$, these two fraction-pairs are equivalent iff
\[
(-a)(-b) = b\cdot a.
\]
Now $(-a)(-b) = ab$ by the product-of-negatives rule for integers, and
$ba = ab$ by commutativity. So the condition holds, and $[(-a,b)] = [(a,-b)]$.

\textbf{Step 7.} Assemble the conclusion.

Combining Steps 5 and 6:
\[
-[(a,b)] = [(-a,b)] = [(a,-b)]. \qedhere
\]
\end{proof}

\begin{remark*}[Proof structure]
The first identity is proved by exhibiting $[(-a,b)]$ as the additive
inverse and invoking uniqueness. The key computation is that the numerator
of the sum $ab + b(-a)$ collapses to $0$. The second identity is a direct
application of the cross-multiplication criterion, using the integer identity
$(-a)(-b) = ab$.
\end{remark*}

\begin{remark*}[Dependencies]~\\
\begin{itemize}
  \item \hyperref[def:rational-operations]{Definition~\ref*{def:rational-operations}}
        \textnormal{(additive inverse must satisfy $[(a,b)] + x = [(0,1)]$)}
  \item \hyperref[def:fraction-pair-equality]{Definition~\ref*{def:fraction-pair-equality}}
  \item \hyperref[thm:ec-addition]{Theorem~\ref*{thm:ec-addition}}
  \item Uniqueness of additive inverses in an abelian group
  \item Basic integer arithmetic: $(-a)(-b) = ab$ and $ab - ab = 0$
\end{itemize}
\end{remark*}
